\documentclass[11pt]{article}
\usepackage[margin=24mm]{geometry}
\usepackage{amsmath,amssymb,amsthm,mathtools,xurl}
\usepackage[hidelinks]{hyperref}
\allowdisplaybreaks\setlength{\emergencystretch}{3em}
\newtheorem{theorem}{Theorem}
\newcommand{\R}{\mathbb R}
\newcommand{\C}{\mathbb C}
\newcommand{\norm}[1]{\left\lVert#1\right\rVert}
\title{Exact finite input reduction for the original quartet-divided moments}
\author{}\date{20 September 2026}
\begin{document}\maketitle

The existing arithmetic-moment provider \cite{AM} proves certified
whole-mass moments of \(g=2\xi\), using the original theta source.
Its accompanying report explicitly leaves division by the quartet
polynomial unevaluated. This note proves the exact finite bridge:
\(4m\) divided even-moment seeds, together with the specified
undivided moments, determine every requested finite divided moment.
It also gives the complete equivalent complex-resolvent input,
all repeated-pole coefficients, and the actual error amplification.
No numerical off-line zero is supplied or asserted.

\section{Original quartet, whole measures, and their exact map}

Retain the programme's actual-quartet domain \cite{MI}:
\begin{equation}\tag{QS1}
\begin{gathered}
 \rho_*=\tfrac12+\delta+i\gamma,\qquad
 0<\delta<\tfrac12,\qquad \gamma>2,\qquad m\in\mathbb Z_{\ge1},\\
 h(s)=\prod_{\epsilon,\eta=\pm1}
       (s-\tfrac12-\epsilon\delta-i\eta\gamma)^m,\qquad g=2\xi .
\end{gathered}
\end{equation}
Here the four selected points are the stipulated actual zero
quartet, with the assigned multiplicity; no existence statement
about an off-line quartet is made. The following identities retain
that same divisor and never replace its parameters by sampled
zero locations.

Define
\begin{equation}\tag{QS2}
\begin{gathered}
 A=2(\delta^2-\gamma^2),\qquad
 B=(\delta^2+\gamma^2)^2,\qquad
 Q(x)=x^2+Ax+B,\\
 a=2m,\quad D=2a=4m,\quad P(x)=Q(x)^a=\sum_{j=0}^D p_jx^j,\\
 u=\gamma^2-\delta^2>0,\qquad v=2\delta\gamma>0,\qquad
 z_\pm=u\pm iv=(\gamma\pm i\delta)^2,\qquad
 R=|z_\pm|=\gamma^2+\delta^2>4.
\end{gathered}
\end{equation}
The symbol \(D\) in this note is the integer degree \(4m\);
the Euler operator used below is denoted \(\mathcal D\).
For the original real coordinate \(y\),
\begin{equation}\tag{QS3}
\begin{aligned}
 h(1/2+iy)
 &=\bigl(((y-\gamma)^2+\delta^2)
          ((y+\gamma)^2+\delta^2)\bigr)^m=Q(y^2)^m,\\
 |h(1/2+iy)|^2&=P(y^2),\\
 Q(x)&=(x-z_+)(x-z_-)=(x-u)^2+v^2\ge v^2
                      \quad(x\ge0).
\end{aligned}
\end{equation}
For the first equality, pairing the two \(\epsilon\) factors
at each fixed \(\eta\) gives
\(-((y-\eta\gamma)^2+\delta^2)\); the two minus signs
cancel before the power \(m\). Expanding their product proves
the displayed \(Q\), including the sign of \(A\).

Use the complete original measures and moments
\begin{equation}\tag{QS4}
\begin{gathered}
 d\mu_y(y)=|g(1/2+iy)|^2\,\frac{dy}{2\pi},\qquad
 d\nu_y(y)=\frac{d\mu_y(y)}{P(y^2)},\\
 \beta_n=\int_{\R}y^{2n}\,d\mu_y(y),\qquad
 \eta_j=\int_{\R}y^j\,d\nu_y(y),\qquad e_n=\eta_{2n}.
\end{gathered}
\end{equation}
Reality and the functional equation of \(g\) imply that both
measures are even. Thus \(\eta_{2n+1}=0\) exactly.
Their masses are \(\beta_0\) and \(e_0\), and no division by
either mass is made.
All moments are finite: \cite{AM}, AM.6--AM.18, proves this
for \(\mu_y\), and QS3 gives \(P(y^2)\ge v^{2a}>0\).

The explicit pushforward \(x=y^2\) is used only to write
the exact polynomial division:
\begin{equation}\tag{QS5}
 d\mu(x)=\frac{|g(1/2+i\sqrt x)|^2}{2\pi\sqrt x}\,dx,\qquad
 d\nu(x)=\frac{d\mu(x)}{P(x)}\quad(x>0),\qquad
 \beta_n=\int_0^\infty x^n\,d\mu,\quad e_n=\int_0^\infty x^n\,d\nu .
\end{equation}
Indeed the two equal halves of the \(y\) integral cancel the
factor \(2\) in \(dy=dx/(2\sqrt x)\).
This proves the exact measure map, with the original \(2\pi\)
and whole mass. It also gives the useful finite guards
\begin{equation}\tag{QS6}
 0<e_j\le v^{-2a}\beta_j,\qquad
 \operatorname{dist}(z_\pm,[0,\infty))=v,\qquad
 |z_+-z_-|=2v=4\delta\gamma .
\end{equation}
Strict positivity follows because the nonzero analytic function
\(g\) vanishes only on a set of real measure zero.

\section{Literal monic recurrence and the entire finite input map}

The coefficients of \(P\), without omitting any term, are
\begin{equation}\tag{QS7}
 p_j=\sum_{\substack{0\le t\le\lfloor j/2\rfloor\\t\ge j-a}}
     \frac{a!}{t!\,(j-2t)!\,(a-j+t)!}\,
          A^{j-2t}B^{a-j+t},\qquad 0\le j\le D.
\end{equation}
This is the multinomial expansion of
\((x^2+Ax+B)^a\). In particular \(p_D=1\),
\(p_{D-1}=aA\), and \(p_0=B^a\).
Multiplying the density relation \(d\mu=P\,d\nu\) by \(x^n\)
and integrating the finite sum proves
\begin{equation}\tag{QS8}
 \beta_n=\sum_{j=0}^Dp_je_{n+j},\qquad
 e_{n+D}=\beta_n-\sum_{j=0}^{D-1}p_je_{n+j}
                  \quad(n\ge0).
\end{equation}
The leading coefficient one is essential; no omitted denominator
or root factor enters this recurrence.

For each \(n\ge0\), let \(q_n,r_n\) be literal Euclidean division:
\begin{equation}\tag{QS9}
 x^n=P(x)q_n(x)+r_n(x),\qquad \deg r_n<D.
\end{equation}
For \(n<D\), \(q_n=0\) and \(r_n=x^n\).
All coefficients can be obtained without root separation
denominators. Define
\begin{equation}\tag{QS10}
\begin{aligned}
 c_\ell
 &=[t^\ell](1+At+Bt^2)^{-a}\\
 &=\sum_{j=0}^{\lfloor\ell/2\rfloor}
   (-1)^{\ell-j}\frac{(a)_{\ell-j}}{(\ell-2j)!\,j!}
                     A^{\ell-2j}B^j,\qquad \ell\ge0 .
\end{aligned}
\end{equation}
Here \((a)_s=a(a+1)\cdots(a+s-1)\) and \((a)_0=1\).
Expanding \((1+z)^{-a}\) formally and then \(z=At+Bt^2\)
proves the coefficient formula. These are finite polynomial
expressions, so no analytic convergence assertion is needed.
For \(n\ge D\), with \(b=n-D\), they give
\begin{equation}\tag{QS11}
\begin{aligned}
 q_n(x)&=\sum_{\ell=0}^{b}c_\ell x^{b-\ell}
         =\sum_{t=0}^{b}q_{n,t}x^t,\qquad q_{n,t}=c_{b-t},\\
 r_n(x)&=\sum_{j=0}^{D-1}r_{n,j}x^j,\qquad
 r_{n,j}=-\sum_{t=0}^{\min(j,b)}p_{j-t}c_{b-t}.
\end{aligned}
\end{equation}
To verify the quotient, multiply
\(x^{-D}\sum_{\ell\ge0}c_\ell x^{-\ell}=1/P(x)\)
as formal Laurent series and keep its nonnegative powers
after multiplication by \(x^n\). All powers at least \(D\)
in \(x^n-Pq_n\) cancel. Its lower coefficients are exactly
the second formula. This proves both QS9 and QS11.

\begin{theorem}[Exact finite seed reduction]
For every \(n\ge0\),
\begin{equation}\tag{QS12}
 e_n=\sum_{j=0}^{D-1}r_{n,j}e_j+
                  \sum_{t=0}^{n-D}q_{n,t}\beta_t ,
\end{equation}
where the second sum is absent if \(n<D\).
Consequently the divided moments through \(\eta_{2L}\),
for a specified integer \(L\ge0\), require only the seeds
\(e_0,\ldots,e_{\min(L,D-1)}\), and, when \(L\ge D\),
the undivided moments \(\beta_0,\ldots,\beta_{L-D}\).
For \(L\ge D-1\) the full seed list is precisely
\(\eta_0,\eta_2,\ldots,\eta_{8m-2}\), with \(4m\) entries.
\end{theorem}
\begin{proof}
Divide QS9 by \(P(x)\) and integrate against the full measure
\(\mu\). Every sum is finite and every moment exists by QS6.
This proves QS12. The degree bounds in QS9 give the stated
input counts. Conversely QS8 recursively determines a unique
sequence from all the \(D\) seeds and the \(\beta_n\).
Thus QS12 and that recurrence give the same exact sequence.
\end{proof}

\section{Exact use of the existing certified arithmetic provider}

The provider \cite{AM} retains
\begin{equation}\tag{QS13}
\begin{gathered}
 \mathcal D=-x\partial_x,\quad
 \phi_*(x)=(4\pi^2x^4-6\pi x^2)e^{-\pi x^2},\quad
 f_0(x)=2\sum_{j\ge1}\phi_*(jx),\\
 H_{\alpha\beta}
 =\langle\mathcal D^\alpha f_0,\mathcal D^\beta f_0\rangle_{L^2(dx)}
 =\int_{\R}\overline{(1/2+iy)^\alpha}(1/2+iy)^\beta\,d\mu_y .
\end{gathered}
\end{equation}
To make the receiving transform completely explicit, put
\[
 T_{n,\alpha}=\binom n\alpha(-1/2)^{n-\alpha}
                       \quad(0\le\alpha\le n).
\]
Binomial expansion of the same original Mellin coordinate gives
\begin{equation}\tag{QS14}
 \beta_n=\sum_{\alpha,\beta=0}^n
        T_{n,\alpha}T_{n,\beta}H_{\alpha\beta}
       =\|(\mathcal D-1/2)^nf_0\|_2^2.
\end{equation}
Indeed the Mellin multiplier in that norm is \((iy)^n\),
whose squared modulus is \(y^{2n}\). Thus this is the exact
AM.17--AM.18 triangular coefficient map, retaining its inverse
\(\binom n\alpha(1/2)^{n-\alpha}\).
For instance \(\beta_0=H_{00}\) is the full arithmetic mass
and \(\beta_1=H_{11}-H_{00}/4\), since
\(H_{01}=H_{00}/2\) by AM.15.

Here are the precise certified-input formulas already proved
in the provider, recorded to specify all finite guards.
Let \(P_0^{\rm AM}(z)=4z^2-6z\),
\(P_{j+1}^{\rm AM}=2z(P_j^{\rm AM}-(P_j^{\rm AM})')\), and
\(Q_j^{\rm AM}=\sum_{\ell=0}^j(-1)^\ell\binom j\ell
P_\ell^{\rm AM}\). Let \(p_{\alpha s}^{\rm AM}\),
\(q_{\alpha s}^{\rm AM}\) be their coefficients and
\(d_\alpha=\alpha+2\). Then AM.8 is the absolutely convergent series
\begin{equation}\tag{QS15}
\begin{aligned}
 H_{\alpha\beta}
 =2\sum_{j,l\ge1}\sum_{s=1}^{d_\alpha}
                    \sum_{t=1}^{d_\beta}
 &(p_{\alpha s}^{\rm AM}p_{\beta t}^{\rm AM}
     +q_{\alpha s}^{\rm AM}q_{\beta t}^{\rm AM})
        \pi^{s+t}j^{2s}l^{2t}\\
 &{}\cdot
 \frac{\Gamma(s+t+1/2,\pi(j^2+l^2))}
      {(\pi(j^2+l^2))^{s+t+1/2}} .
\end{aligned}
\end{equation}
It follows from the complete two-half-domain identity
\[
 H_{\alpha\beta}=\int_1^\infty
 \bigl(f_\alpha f_\beta+\widetilde f_\alpha\widetilde f_\beta\bigr)\,dx,
 \qquad f_\alpha=\mathcal D^\alpha f_0,\quad
 \widetilde f_\alpha=(1-\mathcal D)^\alpha f_0,
 \quad f_\alpha(1/x)=x\widetilde f_\alpha(x).
\]
The source's proof establishes absolute convergence before
interchanging the integral and the double series; both reflected
terms are present in QS15.

For a theta cutoff \(K\ge\max(\alpha,\beta)+2\), put
\begin{equation}\tag{QS16}
\begin{gathered}
 C_\alpha=\sum_s|p_{\alpha s}^{\rm AM}|\pi^s,\qquad
 E_\alpha=\sum_s|q_{\alpha s}^{\rm AM}|\pi^s,\\
 \omega_{\alpha,K}=
       \left(\frac{K+2}{K+1}\right)^{2d_\alpha}e^{-\pi(2K+3)}<1,
 \quad U_{\alpha,K}=\frac{(K+1)^{2d_\alpha}}{1-\omega_{\alpha,K}},\\
 V_{\alpha,K}=\sum_{j=1}^K j^{2d_\alpha}e^{-\pi(j^2-1)}
                 +U_{\alpha,K}e^{-\pi((K+1)^2-1)},\\
 I_s(\lambda)=\frac{\Gamma(s+1/2,\lambda)}
                          {2\lambda^{s+1/2}},\\
 \mathcal E_{\alpha\beta,K}
 =4(C_\alpha C_\beta+E_\alpha E_\beta)
 (U_{\alpha,K}V_{\beta,K}+V_{\alpha,K}U_{\beta,K})
 I_{d_\alpha+d_\beta}(\pi((K+1)^2+1)).
\end{gathered}
\end{equation}
AM.13 proves
\(|H_{\alpha\beta}-H_{\alpha\beta}^{[K]}|
\le\mathcal E_{\alpha\beta,K}\).
Its full tail uses the union \(j>K\) or \(l>K\); no part
of the double complement is discarded.
Thus a certified radius for \(H_{\alpha\beta}\) is its
finite-series arithmetic radius plus QS16, rounded outwards.
The published default provider evaluates degrees \(0\) through
\(5\); higher degrees use the same proved formula and their
stated cutoff guard. No unexecuted higher-degree value is
claimed to have already been numerically certified here.

If entry errors are \(\epsilon_{\alpha\beta}\), QS14 gives
\begin{equation}\tag{QS17}
 |\Delta\beta_n|
 \le\sum_{\alpha,\beta=0}^n
      |T_{n,\alpha}T_{n,\beta}|\epsilon_{\alpha\beta}.
\end{equation}
For a common error \(\epsilon_H\), the bound is
\((3/2)^{2n}\epsilon_H\), since
\(\sum_\alpha|T_{n,\alpha}|=(3/2)^n\).

One may combine QS12 and QS14 \emph{before} enclosing errors.
For \(n\ge D\), set
\begin{equation}\tag{QS18}
 F_{n,\alpha\beta}
 =\sum_{t=\max(\alpha,\beta)}^{n-D}
        q_{n,t}\binom t\alpha\binom t\beta
                     (-1/2)^{2t-\alpha-\beta}.
\end{equation}
Then exactly
\[
 e_n=\sum_{j=0}^{D-1}r_{n,j}e_j+
                  \sum_{\alpha,\beta=0}^{n-D}
                         F_{n,\alpha\beta}H_{\alpha\beta}.
\]
The error multiplier for the \(H\) entries is consequently
\(\sum|F_{n,\alpha\beta}|\epsilon_{\alpha\beta}\), which
retains cancellations among the repeated uses of the same
certified entry. Symmetry permits the exact replacement
of off-diagonal pairs by \(2F_{n,\alpha\beta}H_{\alpha\beta}\).
To receive all \(e_n\), \(n\le L\), a common theta cutoff
\(K\ge L-D+2\) suffices for this source block when \(L\ge D\).

\section{Actual finite error propagation and unavoidable amplification}

For exact \(\delta,\gamma\) and thus exact coefficient polynomials,
write the supplied seed and beta errors as
\(|\Delta e_j|\le\epsilon_j\) and
\(|\Delta\beta_t|\le\epsilon_t^\beta\).
The exact linear map QS12 gives
\begin{equation}\tag{QS19}
 |\Delta e_n|\le
 \sum_{j=0}^{D-1}|r_{n,j}|\epsilon_j+
               \sum_{t=0}^{n-D}|q_{n,t}|\epsilon_t^\beta .
\end{equation}
For a fixed output row this is the exact worst-case
absolute error for independent real input intervals:
choose the signs of all input errors to agree with that
row's real coefficients. Such independently perturbed
inputs need not themselves be the moments of a positive
measure; QS19 describes the information in the separate
input enclosures.

For a requested finite \(L\) define the explicit row sums
\begin{equation}\tag{QS20}
 \kappa_n^{\rm seed}=\sum_{j=0}^{D-1}|r_{n,j}|,\quad
 \kappa_n^\beta=\sum_{t=0}^{n-D}|q_{n,t}|,\quad
 \kappa_L=\max_{0\le n\le L}
                  (\kappa_n^{\rm seed}+\kappa_n^\beta)\ge1 .
\end{equation}
Empty sums are zero. Computing these finite coefficient
polynomials first and supplying every input to absolute
error at most \(\tau/\kappa_L\) proves output error at most
\(\tau\) through \(L\). To reserve \(\tau/2\) for the
outward arithmetic evaluation of the final linear forms,
use input errors at most \(\tau/(2\kappa_L)\).
This is an explicit precision choice, not an assumption
that forward recurrence is stable.

An entirely finite upper majorant is also available.
The factorization in QS2 yields
\[
 c_\ell=\sum_{j=0}^{\ell}
 \binom{a+j-1}{j}\binom{a+\ell-j-1}{\ell-j}
                       z_+^jz_-^{\ell-j},
 \qquad
 |c_\ell|\le\binom{D+\ell-1}{\ell}R^\ell .
\]
The first identity follows by multiplying
\((1-z_+t)^{-a}(1-z_-t)^{-a}\).
Taking moduli and taking the coefficient of \(t^\ell\)
in \((1-t)^{-2a}\) proves the bound.
Let \(W_b=\sum_{\ell=0}^b\binom{D+\ell-1}{\ell}R^\ell\).
Since \(\sum_j|p_j|\le(1+|A|+B)^a\), QS11 implies
\begin{equation}\tag{QS21}
 \kappa_n^\beta\le W_{n-D},\qquad
 \kappa_n^{\rm seed}\le (1+|A|+B)^a W_{n-D}
                                      \quad(n\ge D).
\end{equation}
The computed row sums QS20 usually give a tighter decision.

There is also a necessary lower bound, showing why errors
cannot be assigned by the recurrence order alone.
Because \(x^n-r_n(x)\) is divisible by \((x-z_+)^a\),
its divided derivative of order \(a-1\) vanishes there.
Thus
\[
 \binom n{a-1}z_+^{n-a+1}
 =\sum_{j=a-1}^{D-1}r_{n,j}
                      \binom j{a-1}z_+^{j-a+1}.
\]
For \(n\ge D-1\), taking absolute values and using \(R>1\)
gives
\begin{equation}\tag{QS22}
 \boxed{\quad
 \kappa_n^{\rm seed}\ge
   \frac{\binom n{a-1}}{\binom{D-1}{a-1}}\,
                                  R^{n-D+1}.\quad}
\end{equation}
Indeed \(\binom j{a-1}R^{j-a+1}\) increases on
\(a-1\le j\le D-1\), so its maximum occurs at \(D-1\).
The same bound is attained as an amplification lower bound
by a real homogeneous-error sequence: for a target \(n\),
multiply \(\binom j{a-1}z_+^{j-a+1}\) by a phase making its
\(j=n\) value positive real, take real parts, and divide
by \(\binom{D-1}{a-1}R^{D-a}\).
Its first \(D\) values have modulus at most one. It obeys
the homogeneous QS8 because these divided derivative
sequences are killed by the factor \((E-z_+)^a\) of
\(P(E)\), where \(E\) advances the index. This also follows
directly by differentiating \(z^jP(z)\) \(a-1\) times at \(z_+\).
Its \(n\)th value equals the right side of QS22.

Consequently the precision loss allowed by independent seed
errors includes at least
\((n-D+1)\log_2R+
\log_2\binom n{a-1}-\log_2\binom{D-1}{a-1}\) bits.
This does not assert instability of the underlying positive
integral. It quantifies the cancellation in this particular
exact finite input map.

For uncertain parameter enclosures, preserve correlations
where the arithmetic system can do so; the following bound
is valid without them. Suppose a coefficient \(a_\ell\)
in the finite QS12 or QS18 row is enclosed by
centre \(\widehat a_\ell\), radius \(\rho_\ell\), and its
input \(v_\ell\) by centre \(\widehat v_\ell\), radius
\(\epsilon_\ell\). Then
\begin{equation}\tag{QS23}
 \left|\sum_\ell a_\ell v_\ell-
              \sum_\ell\widehat a_\ell\widehat v_\ell\right|
 \le\sum_\ell\left[
 |\widehat a_\ell|\epsilon_\ell+
                  \rho_\ell(|\widehat v_\ell|+\epsilon_\ell)\right].
\end{equation}
Expand each product as centre plus its two errors to prove it.
The coefficients in QS7, QS10, QS11 and QS18 are finite
polynomials in the original real parameters, so direct
outward polynomial evaluation supplies every \(\rho_\ell\).
Finite parameter boxes used for QS6 must certify
\(\delta_->0\), \(\delta_+<1/2\), and \(\gamma_->2\).
The guard \(v\ge2\delta_-\gamma_->0\) then retains the
original denominator. No guessed zero enclosure is used.

\section{Equivalent finite resolvent input, with all pole coefficients}

The \(D\) real seeds have a second exact presentation.
On \(\C\setminus[0,\infty)\), define
\begin{equation}\tag{QS24}
 \mathcal S(z)=\int_0^\infty\frac{d\mu(x)}{x-z},
 \qquad F_\ell(z)=\int_0^\infty\frac{d\mu(x)}{(x-z)^\ell}
       =\frac{\mathcal S^{(\ell-1)}(z)}{(\ell-1)!}.
\end{equation}
These are holomorphic. On a compact subset of the domain,
distance from the support is positive; the full mass
\(\beta_0\) times that distance to a negative power dominates
every required derivative, proving differentiation under
the integral. Differentiating \((x-z)^{-1}\) produces a
\emph{positive} coefficient at every \(z\) derivative.
Equivalently, with
\(\mathcal R(w)=\int d\mu(x)/(x+w)\),
\begin{equation}\tag{QS25}
 F_\ell(z)=\frac{(-1)^{\ell-1}}{(\ell-1)!}
                \mathcal R^{(\ell-1)}(-z).
\end{equation}
This records the different sign of differentiation in the
plus-denominator parameter. In either convention,
\(|F_\ell(z_\pm)|\le\beta_0/v^\ell\).

For every integer \(n\ge0\), every \(1\le\ell\le a\), and
\(\sigma\in\{+,-\}\), define
\begin{equation}\tag{QS26}
 C_{\sigma,\ell}^{(n)}
 =\sum_{t=0}^{\min(n,a-\ell)}
  \binom nt z_\sigma^{n-t}(-1)^{a-\ell-t}
  \binom{2a-\ell-t-1}{a-\ell-t}
      (z_\sigma-z_{-\sigma})^{-(2a-\ell-t)} .
\end{equation}
Then the complete partial-fraction formula is
\begin{equation}\tag{QS27}
 \frac{x^n}{P(x)}
 =q_n(x)+\sum_{\sigma=\pm}\sum_{\ell=1}^a
                      \frac{C_{\sigma,\ell}^{(n)}}{(x-z_\sigma)^\ell}.
\end{equation}
To prove the coefficient, expand
\(x^n(x-z_{-\sigma})^{-a}\) at \(x=z_\sigma\).
The numerator has Taylor coefficient
\(\binom nt z_\sigma^{n-t}\);
the other factor's coefficient of order \(s\) is
\((-1)^s\binom{a+s-1}{s}
(z_\sigma-z_{-\sigma})^{-a-s}\).
Taking \(s+t=a-\ell\) gives QS26.
Subtracting these complete principal parts from \(x^n/P\)
leaves its polynomial quotient \(q_n\): the proper remainder
has no pole and tends to zero at infinity, and is therefore zero.
This proves every repeated-pole coefficient with its sign.

Since the parameters and measure are real,
\(C_{-,\ell}^{(n)}=\overline{C_{+,\ell}^{(n)}}\) and
\(F_\ell(z_-)=\overline{F_\ell(z_+)}\). Integration yields
\begin{equation}\tag{QS28}
 e_n=\sum_{t=0}^{n-D}q_{n,t}\beta_t+
          2\Re\sum_{\ell=1}^a C_{+,\ell}^{(n)}F_\ell(z_+).
\end{equation}
In particular all \(4m=D\) real seeds are obtained from the
\(a=2m\) complex values \(F_1(z_+),\ldots,F_a(z_+)\).
The exact reverse map is
\begin{equation}\tag{QS29}
\begin{aligned}
 F_\ell(z_+)&=\sum_{j=0}^{D-\ell}d_{\ell,j}e_j,\\
 d_{\ell,j}
 &=\sum_{\substack{u+w=j\\0\le u\le a-\ell\\0\le w\le a}}
    \binom{a-\ell}{u}\binom aw
       (-z_+)^{a-\ell-u}(-z_-)^{a-w}.
\end{aligned}
\end{equation}
Indeed the numerator polynomial in
\[
 \frac1{(x-z_+)^\ell}
   =\frac{(x-z_+)^{a-\ell}(x-z_-)^a}{P(x)}
\]
has exactly these coefficients and degree \(D-\ell<D\).
Partial-fraction uniqueness proves that these two real-linear
maps are inverse: the \(D\) functions \(x^j/P\), \(0\le j<D\),
and the \(a\) conjugate pairs \((x-z_\pm)^{-\ell}\) are
two bases of the same \(D\)-dimensional space of real
proper rational functions. Thus the resolvent presentation
retains precisely all seed data rather than assuming extra
moments can be recovered from the \(\beta_n\) alone.

For certified disks of radii \(\epsilon_\ell^F\) for these
resolvent values, QS28 gives the explicit bound
\begin{equation}\tag{QS30}
 |\Delta e_n|\le
 \sum_{t=0}^{n-D}|q_{n,t}|\epsilon_t^\beta+
                    2\sum_{\ell=1}^a
                            |C_{+,\ell}^{(n)}|\epsilon_\ell^F .
\end{equation}
Use QS23 as well when the pole coefficients are enclosed.
All denominators in QS26 have the positive guard
\(|z_+-z_-|\ge4\delta_-\gamma_-\).
The factors in QS26 can grow when \(\delta\) is small;
the polynomial formulas QS10--QS12 supply an exact alternative
without those root-separation denominators. Neither
representation is assigned stability without its actual
finite row sums.

\section{Exact receiver for the original tensor moments}

For \(k\ge1\), keep \(m_k=w_h^{*k}\), where
\(w_h(y)\,dy=d\nu_y(y)\), and its original sum coordinate.
Let \(\zeta_j=\int_{\R}y^jm_k(y)\,dy\).
For any finite degree bound \(2L\), expansion of
\((y_1+\cdots+y_k)^j\), justified by the finite absolute
moments already proved, gives
\begin{equation}\tag{QS31}
 \zeta_j=j![z^j]\left(\sum_{n=0}^L
                     \frac{e_nz^{2n}}{(2n)!}\right)^k
                 \quad(0\le j\le2L).
\end{equation}
In particular \(\zeta_0=e_0^k\), and all odd moments are
zero. This is the original finite convolution map in
\cite{MI}, with the entire arithmetic mass retained.
If the input centres are \(\widehat e_n\), with radii
\(\epsilon_n\), set the finite coefficient polynomials
\[
 E(z)=\sum_{n=0}^L\frac{|\widehat e_n|z^{2n}}{(2n)!},
 \qquad
 E_+(z)=\sum_{n=0}^L
             \frac{(|\widehat e_n|+\epsilon_n)z^{2n}}{(2n)!}.
\]
Then
\begin{equation}\tag{QS32}
 |\Delta\zeta_j|\le j![z^j]\bigl(E_+(z)^k-E(z)^k\bigr).
\end{equation}
Expand each product into centres and errors. All terms
containing at least one error are bounded by the corresponding
term in this nonnegative coefficient difference. This proves
QS32 without replacing the original source by a probability
measure or by an average-coordinate tensor.

\begin{figure}[ht]\centering
\(\begin{array}{ccc}
 (e_0,\ldots,e_{D-1})&
 \longleftrightarrow&
 (F_1(z_+),\ldots,F_a(z_+))\\
 \scriptstyle{\rm QS12}\downarrow&&
 \downarrow\scriptstyle{\rm QS28}\\
 (e_0,\ldots,e_L)&=&(e_0,\ldots,e_L)\\
 &\downarrow\scriptstyle{\rm QS31}&\\
 &(\zeta_0,\ldots,\zeta_{2L})&
\end{array}\)
\caption{The exact seed and resolvent presentations.
Each downward seed reduction also receives the existing
undivided moments \(\beta_0,\ldots,\beta_{L-D}\) when
\(L\ge D\), through the original AM coefficient map
QS13--QS18. The top arrow and its inverse are QS28--QS29.
The bottom arrow retains total mass \(e_0^k\).
The errors of these maps are QS19, QS23, QS30 and QS32.
No seed centre is supplied by the diagram itself.}
\end{figure}

\section{What is proved and checked}

The finite source requirement is now exact: at most \(4m\)
real divided seeds, or equivalently \(2m\) complex resolvent
jets, plus the specified already available arithmetic-moment
provider. QS19--QS23 make every required finite input accuracy
computable; QS22 proves a real worst-case amplification that
a numerical implementation must retain. These results do not
claim numerical values for the hypothetical quartet's seeds.
Their certified original integral evaluation is the separate
receiving calculation.

The accompanying exact checker verifies coefficient formulas,
Euclidean division, the monic recurrence, repeated-pole
coefficients, both Stieltjes derivative signs, both inverse
seed maps, the AM triangular receiver, full-mass convolution,
and the error identities on finite rational measures.
Its rational \(\delta,\gamma\) choices are algebraic test
parameters, explicitly not claimed zeta zeros. Actual programme
use keeps QS1 and every original coefficient.

\begin{thebibliography}{9}
\bibitem{AM} Split-Zero research programme,
\emph{Certified moments from the original arithmetic theta source},
12 September 2026,
\path{work/cumulative_deligne_build_20260913_v10/tex/arithmetic_moments.tex},
AM.1--AM.19, particularly the full double series AM.8,
the complete omitted-tail bound AM.13, and the original
triangular coefficient map AM.17--AM.18.
The complete proof provider was read for this derivation.
Its companion report is
\texttt{work/arithmetic\_moment\_exact\_20260912.md}.
The published executed finite certificates cover the degrees
stated in that report; the general formulas are the provider's
proved all-degree theorem.
\bibitem{MI} Split-Zero research programme,
\emph{Moving-interval native margins and finite raw-moment precision},
20 September 2026,
\texttt{MOVING\_INTERVAL\_NATIVE\_PRECISION.tex},
MI1 for the actual full-quartet domain and MI14--MI17
for the finite convolution receiver. The present QS7--QS30
supplies the exact input reduction preceding that receiver.
\end{thebibliography}
\end{document}
